% Вторая часть раздела «Автономные беспилотные миссии»
\subsection{Проверка знаний}
\begin{enumerate}
  \item Как инициировать предстартовую подготовку? \texttt{ap.push(\_\_\_)} 
  \begin{enumerate}[label=(\alph*)] \item \texttt{Ev.MCE\_PREFLIGHT} \item \texttt{Ev.PREFLIGHT} \item \texttt{Ev.START} \item \texttt{Ev.PRE} \end{enumerate}
  Ответ: (a). Пояснение: используется пространство имён \texttt{Ev.MCE\_*}.
  \item Чем реагировать на завершение взлёта? 
  \begin{enumerate}[label=(\alph*)] \item \texttt{callback(Ev.TAKEOFF\_COMPLETE)} \item \texttt{function callback(event) if event==Ev.TAKEOFF\_COMPLETE then ... end end} \item \texttt{ap.onTakeoff()} \item \texttt{Timer.new(5, ...)} \end{enumerate}
  Ответ: (b). Пояснение: события обрабатываются внутри \texttt{callback}.
  \item Как задать полёт к точке? 
  \begin{enumerate}[label=(\alph*)] \item \texttt{ap.goToLocalPoint(x,y,z)} \item \texttt{ap.fly(x,y,z)} \item \texttt{ap.moveTo(x,y,z)} \item \texttt{ap.goto(x,y,z)} \end{enumerate}
  Ответ: (a). Пояснение: корректный API платформы.
  \item Как обновить курс по градусам \texttt{deg}? 
  \begin{enumerate}[label=(\alph*)] \item \texttt{ap.updateYaw(deg)} \item \texttt{ap.updateYaw(deg*math.pi/180)} \item \texttt{ap.updateYaw(deg/180)} \item \texttt{ap.updateYaw(math.pi*deg)} \end{enumerate}
  Ответ: (b). Пояснение: \texttt{updateYaw} принимает радианы.
  \item Как безопасно выполнить действие через 2 секунды? 
  \begin{enumerate}[label=(\alph*)] \item \texttt{sleep(2); action()} \item \texttt{Timer.callLater(2, action)} \item \texttt{Timer.new(2, action)} \item \texttt{ap.delay(2, action)} \end{enumerate}
  Ответ: (b). Пояснение: отложенный неблокирующий вызов.
  \item Как задать таблицу точек окружности? 
  \begin{enumerate}[label=(\alph*)] \item \texttt{{cos(i),sin(i)}} \item \texttt{{r*math.cos(...), r*math.sin(...)}} \item \texttt{{r*cos(...), r*sin(...)}} \item \texttt{{x=i, y=i}} \end{enumerate}
  Ответ: (b). Пояснение: вычисление по параметрическим формулам.
  \item Как корректно выключить двигатели? 
  \begin{enumerate}[label=(\alph*)] \item \texttt{ap.push(ENGINES\_DISARM)} \item \texttt{ap.push(Ev.ENGINES\_DISARM)} \item \texttt{ap.disarm()} \item \texttt{ap.push(Ev.MCE\_STOP)} \end{enumerate}
  Ответ: (b). Пояснение: требуется пространство имён \texttt{Ev}.
  \item Как реагировать на удар/столкновение? 
  \begin{enumerate}[label=(\alph*)] \item \texttt{event==Ev.SHOCK} \item \texttt{event==Ev.COLLISION} \item \texttt{event==Ev.CRASH} \item \texttt{event==Ev.HIT} \end{enumerate}
  Ответ: (a). Пояснение: используйте \texttt{Ev.SHOCK}.
  \item Как организовать непрерывное обновление точки окружности? 
  \begin{enumerate}[label=(\alph*)] \item \texttt{while ... do ap.goToLocalPoint(...) end} \item \texttt{Timer.new(period, function() ap.goToLocalPoint(...) end)} \item \texttt{repeat ap.goToLocalPoint(...) until false} \item \texttt{sleep(0.1); ap.goToLocalPoint(...)} \end{enumerate}
  Ответ: (b). Пояснение: периодический таймер без блокировок.
  \item Как прочитать RC-каналы? 
  \begin{enumerate}[label=(\alph*)] \item \texttt{Sensors.rc()} \item \texttt{ap.rc()} \item \texttt{rc.read()} \item \texttt{Sensors.readRC()} \end{enumerate}
  Ответ: (a). Пояснение: API сенсоров.
  \item Как проверить положение тумблера SwA в \texttt{rc\_chans}? 
  \begin{enumerate}[label=(\alph*)] \item \texttt{rc\_chans[1]} \item \texttt{rc\_chans[8]} \item \texttt{rc\_chans["SwA"]} \item \texttt{rc\_chans[0]} \end{enumerate}
  Ответ: (b). Пояснение: восьмой канал.
  \item Как построить переходы состояний? 
  \begin{enumerate}[label=(\alph*)] \item \texttt{if ... then ...} \item таблица функций \texttt{action[state]()} \item \texttt{switch(state)} \item \texttt{goto state} \end{enumerate}
  Ответ: (b). Пояснение: идиоматичный подход через таблицу функций.
  \item Как задать высоту полёта по окружности? 
  \begin{enumerate}[label=(\alph*)] \item \texttt{ap.setHeight(z)} \item \texttt{ap.goToLocalPoint(x,y,z)} \item \texttt{ap.updateYaw(z)} \item \texttt{Sensors.altimeter(z)} \end{enumerate}
  Ответ: (b). Пояснение: высота \texttt{z} задаётся в целевой точке.
  \item Как обеспечить аварийную индикацию при низком напряжении? 
  \begin{enumerate}[label=(\alph*)] \item в \texttt{callback} по \texttt{Ev.LOW\_VOLTAGE2} остановить таймеры и включить красный \item задержка через \texttt{sleep} \item увеличить \texttt{height} \item проигнорировать событие \end{enumerate}
  Ответ: (a). Пояснение: минимизация нагрузки и чёткая реакция.
\end{enumerate}

\subsection*{Итоги раздела}
Вы научились использовать события платформы, интерфейсы \texttt{ap.*}, таймеры и индикаторы для построения автономных миссий: от окружности и многоугольника до RC‑управления и аварийных реакций. Практические задания позволяют закрепить навыки и подготовить собственные сценарии полётов.

\subsection{Задания для лабораторной работы}
\subsubsection{Базовые миссии}
\begin{enumerate}
  \item \textbf{Полёт по окружности радиуса \texttt{r=0.5} на высоте \texttt{z=1.0}.}
    Требуется реализовать алгоритм, обеспечивающий устойчивый полёт по окружности: выполнить \texttt{ap.push(Ev.MCE\_PREFLIGHT)} и \texttt{ap.push(Ev.MCE\_TAKEOFF)}, запускать обновление точки через \texttt{Timer.new(0.1, fn)} после \texttt{Ev.TAKEOFF\_COMPLETE}, вычислять \texttt{x=r*cos(theta)}, \texttt{y=r*sin(theta)} с переводом градусов в радианы; при необходимости сопровождать курс вызовом \texttt{ap.updateYaw(angle*math.pi/180)}; использовать индикацию \texttt{Ledbar.new(25)} для этапов; по \texttt{Ev.SHOCK} — останов таймеров и \texttt{ap.push(Ev.ENGINES\_DISARM)}, по \texttt{Ev.LOW\_VOLTAGE2} — посадку.
    \textit{Ожидаемый результат:} квадрокоптер следует окружности радиуса \texttt{0.5} на высоте \texttt{1.0}, выполняет посадку, индикация отражает этапы.
  \item \textbf{Полёт по многоугольнику из \texttt{N} точек окружности.}
    Требуется реализовать генерацию точек по равномерным углам \texttt{a=360*i/N}, переводить углы в радианы при вычислении координат, стартовать маршрут по \texttt{Ev.TAKEOFF\_COMPLETE} и переходить к следующей точке по \texttt{Ev.POINT\_REACHED}; обеспечить безопасные реакции и индикацию как в окружности.
    \textit{Ожидаемый результат:} последовательный обход всех \texttt{N} точек и посадка после завершения маршрута.
  \item \textbf{Траектория «восьмёрка».}
    Требуется реализовать два состояния (\texttt{FIRST}, \texttt{SECOND}) для обхода двух окружностей с противоположным направлением, сместить центры по оси X, выбрать шаг угла разного знака для фаз, завершить миссию возвратом в \texttt{(0,0,z)} и посадкой.
    \textit{Ожидаемый результат:} полёт по «восьмёрке» без рывков, завершение в центре и посадка.
  \item \textbf{Спираль подъёма.}
    Требуется реализовать параметрическую спираль с приращениями \texttt{r=r+dr} и \texttt{z=z+dz} на каждом тике; по достижении \texttt{z\_max} остановить таймер, перейти в безопасную точку базы и выполнить посадку.
    \textit{Ожидаемый результат:} плавный подъём до \texttt{z\_max}, возврат к базе и посадка.
  \item \textbf{Сопровождение курса \texttt{yaw}.}
    Требуется синхронизировать \texttt{ap.updateYaw} с параметром окружности, использовать радианы; ограничить шаг изменения угла разумным значением, чтобы исключить резкие скачки.
    \textit{Ожидаемый результат:} курс следует траектории без резких изменений, ориентация корректна.
  \item \textbf{Запуск и останов миссии RC-тумблером.}
    Требуется считывать \texttt{Sensors.rc()}, анализировать \texttt{rc\_chans[8]}; при включении тумблера выполнять предстарт и взлёт, при выключении — останавливаться, переходить в \texttt{(0,0,z)} и выполнять посадку; опрашивать RC через периодический \texttt{Timer.new(0.5, fn)}.
    \textit{Ожидаемый результат:} надёжный старт/останов миссии по положению тумблера с безопасным завершением.
  \item \textbf{Аварийная реакция на \texttt{Ev.SHOCK}.}
    Требуется немедленно остановить все таймеры миссии, включить красный цвет на всех \texttt{0..24} индексах линейки и выполнить \texttt{ap.push(Ev.ENGINES\_DISARM)}.
    \textit{Ожидаемый результат:} мгновенное выключение двигателей и красная индикация аварии.
  \item \textbf{Посадка при низком напряжении \texttt{Ev.LOW\_VOLTAGE2}.}
    Требуется остановить миссию, перейти в безопасную точку и выполнить \texttt{ap.push(Ev.MCE\_LANDING)}, включить красную индикацию низкого напряжения.
    \textit{Ожидаемый результат:} безопасная посадка при низком напряжении, отчёт индикацией.
  \item \textbf{Обход квадрата по точкам, два цикла.}
    Требуется задать вершины \texttt{(-a/2,-a/2,z)}, \texttt{(a/2,-a/2,z)}, \texttt{(a/2,a/2,z)}, \texttt{(-a/2,a/2,z)}, выполнять переходы по \texttt{Ev.POINT\_REACHED} и совершить два полных обхода перед посадкой.
    \textit{Ожидаемый результат:} два полных обхода квадрата и посадка в конце маршрута.
  \item \textbf{«Лоитер» по малой окружности \texttt{r=0.2}.}
    Требуется периодически обновлять точку окружности таймером, обеспечить безопасную остановку по событию или RC, сопровождать миссию световой индикацией этапов.
    \textit{Ожидаемый результат:} стабильное удержание по малой окружности и корректная остановка по команде.
\end{enumerate}

\subsubsection{Продвинутые задания}
\begin{enumerate}
  \item \textbf{Плавная смена цвета этапов миссии.}
    Требуется сопоставить цвета состояниям миссии (подготовка/полёт/посадка/ожидание), обновлять все индексы \texttt{0..24} без блокировок \texttt{sleep()}, реализовать смену через функции стадий и отложенные вызовы.
    \textit{Ожидаемый результат:} плавная индикация переходов между этапами без пропусков и блокировок.
  \item \textbf{Динамический радиус окружности по RC7.}
    Требуется считывать \texttt{Sensors.rc()}, нормировать значение канала к \texttt{[0,1]}, вычислять радиус \texttt{r=0.2+0.8*norm} и периодически обновлять точку окружности.
    \textit{Ожидаемый результат:} радиус плавно меняется от \texttt{0.2} до \texttt{1.0} по RC7, траектория стабильна.
  \item \textbf{Ограничение скорости изменения угла \texttt{angle}.}
    Требуется реализовать \texttt{advanceAngle(step)} с ограничением \texttt{step} диапазоном \texttt{[-maxStep,maxStep]} и применять его при каждом обновлении параметра окружности.
    \textit{Ожидаемый результат:} угол изменяется с ограниченной скоростью, резкие скачки отсутствуют.
  \item \textbf{Перезапуск миссии после \texttt{Ev.COPTER\_LANDED}.}
    Требуется по событию выполнить \texttt{Timer.callLater(3, fn)} и инициировать новый взлёт, исключая блокирующие задержки.
    \textit{Ожидаемый результат:} автоматический повтор миссии после посадки с паузой, без блокировок.
  \item \textbf{«Морзянка» этапов на линейке.}
    Требуется реализовать мигание заданное число раз для каждого этапа, используя короткие функции и отложенные вызовы.
    \textit{Ожидаемый результат:} корректный световой отчёт по морзянке для каждого этапа.
  \item \textbf{Траектория «лепесток».}
    Требуется реализовать формулу \texttt{r=r0*(1+0.2*sin(k*theta))} при \texttt{theta=angle*pi/180}, выбрать \texttt{k>=3} и обновлять точку параметрически.
    \textit{Ожидаемый результат:} визуально различимый «лепесток» по параметрической траектории.
  \item \textbf{Проверка точности навигации.}
    Требуется обход 8 точек окружности, измерение интервалов \texttt{Δt} между \texttt{Ev.POINT\_REACHED} с использованием \texttt{os.clock()}.
    \textit{Ожидаемый результат:} получены интервалы времени, оценено постоянство периода между достижениями точек.
  \item \textbf{Миссия «периметр» вокруг базы.}
    Требуется определить набор точек по квадрату вокруг \texttt{(0,0)}, включить возврат в исходную точку и выполнить посадку.
    \textit{Ожидаемый результат:} полный обход периметра и посадка в исходной точке.
  \item \textbf{«Окно безопасности» для \texttt{yaw}.}
    Требуется ограничить изменение \texttt{yaw} на тик не более \texttt{maxYawRate*dt} и обновлять точку только при соблюдении ограничения.
    \textit{Ожидаемый результат:} курс изменяется не быстрее установленной скорости, миссия сохраняет устойчивость.
  \item \textbf{Демонстрационный сценарий (RC старт, \(\approx\) 2 минуты, посадка).}
    Требуется запуск по RC, полёт по окружности около двух минут, посадка, световой отчёт этапов, отсутствие блокировок и проверка безопасного завершения.
    \textit{Ожидаемый результат:} демонстрационный полёт по окружности, корректная посадка и полная индикация этапов.
\end{enumerate}

% Итоговые выводы опущены ради упрощения объёма задания.
